\documentclass[11pt,a4paper]{article}
\usepackage[margin=25mm]{geometry}
\usepackage{amsmath,amssymb,amsthm}
\usepackage[hidelinks]{hyperref}
\newtheorem{theorem}{Theorem}
\newtheorem{lemma}{Lemma}
\newcommand{\Z}{\mathbb Z}
\newcommand{\wt}{\operatorname{wt}}
\newcommand{\eps}{\varepsilon}
\title{A nine-vertex counterexample to the crystal duality criterion}
\author{Source 00000003849: local formalization package}
\date{8 October 2026}
\begin{document}
\maketitle

\noindent\textbf{Claim addressed.} The source asserts that a crystal is
isomorphic to its arrow-reversal dual if and only if its weight set is centrally
symmetric, together with a polynomial-time verification claim. The finite
counterexample below refutes the sufficient direction. Both source languages
say weight \emph{set}, with no connectedness or irreducibility restriction.
No complexity theorem is established here.

\section*{The root datum and crystal convention}
Use the type-$A_2$ weight lattice $P=\Z^2$ in fundamental-weight coordinates.
The simple roots and their coroot pairings are
\[
 \alpha_1=(2,-1),\qquad \alpha_2=(-1,2),\qquad
 \langle(a,b),\alpha_1^\vee\rangle=a,\quad
 \langle(a,b),\alpha_2^\vee\rangle=b.
\]
Thus the Cartan matrix has diagonal entries $2$ and off-diagonal entries $-1$.
The six roots are $\pm\alpha_1$, $\pm\alpha_2$ and
$\pm(\alpha_1+\alpha_2)$, with corresponding coroots
$\pm(1,0)$, $\pm(0,1)$ and $\pm(1,1)$. The pairing of the weight and coweight
lattices is the standard dot product. Lean checks that corresponding roots and
coroots pair to $2$, are distinct, and are closed under the corresponding
simultaneous reflections. In particular, the simple reflections are
\[
 s_1(a,b)=(-a,a+b),\qquad s_2(a,b)=(a+b,-b).
\]

We use the finite-type Kashiwara crystal convention. A crystal consists of a
weight map, partial operators $e_i,f_i$ and integer string functions
$\eps_i,\varphi_i$. The partial inverse relation and lowering relations are
\begin{align*}
 f_i(x)=y&\iff e_i(y)=x,\\
 \wt(y)&=\wt(x)-\alpha_i, &
 \eps_i(y)&=\eps_i(x)+1, &
 \varphi_i(y)&=\varphi_i(x)-1,
\end{align*}
with the corresponding raising relations and
$\varphi_i(x)=\eps_i(x)+\langle\wt(x),\alpha_i^\vee\rangle$.
The absent vertex is Lean's \texttt{Option.none}. These are a specialization of
the usual abstract axioms to finite string values, not assumed admissibility.
All axioms are proved for the constructed crystal. Exact nonzero operator
string lengths are also proved, so the example is seminormal.

\section*{The actual finite crystal}
Let $S$ have three vertices $p_0,p_1,p_2$, weights
\[
 (1,0),\qquad (-1,1),\qquad (0,-1),
\]
and arrows $p_0\xrightarrow{1}p_1\xrightarrow{2}p_2$.
For each color, the raising operator is the reverse partial arrow; its string
functions are $0$ or $1$ according to whether the corresponding arrow exists.
This is the standard fundamental $A_2$ crystal $B(\omega_1)$.
For an explicit weight correspondence, the standard $GL_3$ weights $e_1,e_2,e_3$
map under
\[
 (u_1,u_2,u_3)\longmapsto (u_1-u_2,u_2-u_3)
\]
to the three listed weights. The same map sends the two $GL_3$ simple roots
to $\alpha_1,\alpha_2$, and kills the central direction. These identities are
also proved in Lean. The name of the component is therefore supported by its
precise graph and weight correspondence; quantum-group representation theory
is not needed as an unproved premise.

For any crystal $C$, its dual $C^\vee$ has
\[
 \wt^\vee_C(v)=-\wt_C(v),\quad
 e_i^\vee=f_i^C,\quad f_i^\vee=e_i^C,\quad
 \eps_i^\vee=\varphi_i^C,\quad
 \varphi_i^\vee=\eps_i^C.
\]
The generic Lean construction proves all crystal axioms for this dual.
The dual of $S$ has highest weight $(0,1)$ and is the standard dual component
$B(\omega_2)$.

Set $B=S\sqcup S\sqcup S^\vee$. Its vertex set is
$\{0,1,2\}\times\{0,1,2\}$. The first coordinate selects the component.
All operators preserve that coordinate. The two first components use $S$;
the last uses the actual constructed dual. Lean proves the component
correspondence for weights, both operators and both string functions.

\begin{center}
\begin{tabular}{ccc}
\hline
Component & Weights at positions $0,1,2$ & Arrows\\
\hline
$0$ & $(1,0),(-1,1),(0,-1)$ & $0\xrightarrow{1}1\xrightarrow{2}2$\\
$1$ & $(1,0),(-1,1),(0,-1)$ & $0\xrightarrow{1}1\xrightarrow{2}2$\\
$2$ & $(-1,0),(1,-1),(0,1)$ & $2\xrightarrow{2}1\xrightarrow{1}0$\\
\hline
\end{tabular}
\end{center}

\section*{Symmetry and non-isomorphism}
The weight \emph{set} of $B$ is
\[
 \{\pm(1,0),\;\pm(-1,1),\;\pm(0,-1)\}.
\]
It is invariant under negation, hence centrally symmetric with centre zero.
Lean states this for the actual set $\operatorname{range}(\wt_B)$, using the
definition $\exists c\in P\;\forall w\in P:\;w\in X\iff 2c-w\in X$.

\begin{lemma}
No bijection $g:B\to B$ satisfies $\wt_B(g(v))=-\wt_B(v)$ for every vertex.
\end{lemma}
\begin{proof}
The distinct vertices $(0,0)$ and $(1,0)$ both have weight $(1,0)$.
The only vertex of weight $(-1,0)$ is $(2,0)$. The required condition would
send both distinct vertices to $(2,0)$, contradicting injectivity.
\end{proof}

\begin{theorem}
The admissible nine-vertex crystal $B$ has centrally symmetric weight set, but
$B^\vee$ is not isomorphic to $B$ as a crystal.
\end{theorem}
\begin{proof}
A crystal isomorphism is an actual bijection preserving weights, both crystal
operators and both string functions. If $g:B^\vee\to B$ were such an
isomorphism, weight preservation and the definition of the dual would give
$\wt_B(g(v))=\wt_B^\vee(v)=-\wt_B(v)$, contrary to the lemma.
The reverse isomorphism direction is excluded by the same proved bridge.
\end{proof}

This refutes the universal sufficient direction and the source's iff already
within nonempty finite-type $A_2$ crystals. A counterexample in that subclass
also refutes an assertion made for unqualified crystals. It does not replace
the conjecture by a connected-crystal assertion or a multiset assertion.
The algorithmic conjunct need not be resolved to refute a conjunction whose
iff conjunct is false.

\section*{Formal artifact correspondence}
The Main module imports the formal crystal source. In the namespace
\texttt{CrystalCounterexample}, the key declarations are listed below.
\begin{quote}\ttfamily\small
package/CrystalCounterexample.lean\\
original\_sufficient\_direction\_false\\
original\_iff\_false\\
Crystal; Crystal.dual; Crystal.Iso\\
Witness.crystal; Witness.seminormal\\
witness\_centrallySymmetric; Witness.not\_selfDual
\end{quote}

Preparation compiles use official Lean 4.33.0 and fixed mathlib revision
\begin{quote}\ttfamily\small
db584cd6d46c92f209a44c0f1c829460d327499d
\end{quote}
Printed transitive axiom dependencies are limited to the standard three:
\begin{quote}\ttfamily\small
propext; Classical.choice; Quot.sound
\end{quote}
Preparation evidence is not manager acceptance:
independent source-semantic review and a clean deterministic gate remain
separate steps. No polynomial-time result is claimed.

\begin{thebibliography}{1}
\bibitem{marberg}
E. Marberg, \emph{Combinatorics of crystal bases}, MATH 6150I, Lecture 5,
HKUST, Spring 2020. Definitions 2.1 and 3.4 and Section 4 supply the standard
crystal, dual and isomorphism conventions.
\url{https://www.math.hkust.edu.hk/~emarberg/teaching/2020/Math6150I/lectures/05_Math6150I_Spring2020.pdf}
\end{thebibliography}
\end{document}
